\begin{claim}{C.7b}{Proved}
Let $p$ be any prime and $K/\Q_p$ finite, put $F_K=K\Q_p(\mu_\infty)$,
and let $V/K$
be a hyperbolic curve of genus zero with $x\in V(K)$. Every element of
$\pi_1(V_{\bar K})$ fixed by $s_x(G_{F_K})$ has zero image in its
\emph{full} geometric abelianization.

For the explicit tripod over $\Q_5$, let $\Pi$ be its full geometric
fundamental group, choose cusp generators $a,b$, and put
\[
 \mathcal N=\ker(\Pi\longrightarrow\widehat\Z),\qquad
 a\longmapsto1,\quad b\longmapsto0,
 \qquad
 \mathcal Q=\Pi/\overline{[\mathcal N,\mathcal N]}.
\]
Then
\[
 \operatorname{im}\bigl(\Pi^{s_2(H)}\to\mathcal Q\bigr)=1,
 \qquad
 \operatorname{im}\bigl(\Pi^{s_{\overrightarrow{01}}(H)}
                  \to\mathcal Q\bigr)=\overline{\langle a\rangle}
                  \simeq\widehat\Z,
 \qquad H=G_{\Q_5(\mu_\infty)}.
\]
In particular the tangent fixed group's full abelianization image is
exactly the cusp-$0$ line. These assertions determine images of full fixed
elements, and do not assert that the residual nonabelian kernels vanish.
\end{claim}
\begin{proof}
The elementary restriction argument (C.7a.1) extends to \emph{every}
finite extension $M/\Q_p$, without requiring $M\subset\Q_p(\mu_\infty)$.
Indeed
\[
 \operatorname{Gal}\bigl(M\Q_p(\mu_\infty)/M\bigr)
\]
is abelian, as a subgroup of the Galois group of the cyclotomic extension
of $\Q_p$, and its $p$-adic cyclotomic character has infinite image.
The same cocycle identity proves
\[
 H^1(G_M,\Q_p(1))\ \hookrightarrow\
 H^1\bigl(G_{M\Q_p(\mu_\infty)},\Q_p(1)\bigr).               \tag{C.7b.1}
\]
Consequently, if $u\in M^\times$ has zero Kummer class on the group on
the right, then $u$ is a root of unity. Its valuation must be zero;
on the unit group the kernel of the local logarithm is exactly its finite
torsion subgroup, by the logarithm isomorphism on an open principal-unit
subgroup. This conclusion only concerns elements of the \emph{finite}
local field $M$, and does not assert that the cyclotomic extension has
finite torsion.

We first treat $V=\mathbf P^1\setminus\{0,1,\infty\}$ and an arbitrary
$t\in K\setminus\{0,1\}$. Let $g$ be a full geometric element fixed by
$G_{F_K}$. For any $m\ge3$ use the finite etale cyclic cover
\[
 V_m=\mathbf P^1\setminus\bigl(\{0,\infty\}\cup\mu_m\bigr)
       \longrightarrow V,\qquad x=z^m.
\]
Choose $\alpha$ with $\alpha^m=t$. The lift of the loop $g$ starting
at $\alpha$ ends at $\zeta\alpha$ for some $\zeta\in\mu_m$; this $\zeta$
is precisely its winding coordinate at $0$ modulo $m$. Choose a finite
extension $M/K$ containing $\alpha$ and $\mu_m$, and restrict to
$G_{M\Q_p(\mu_\infty)}$. All endpoints and cusps are then rational,
and the lifted etale path is Galois invariant: its image $g$ is invariant
and its starting point is fixed, so uniqueness of path lifting proves the
assertion. For every $\eta\in\mu_m$, the invertible function $z-\eta$
therefore gives a Galois-invariant path on $\mathbf G_m$ from
$\alpha-\eta$ to $\zeta\alpha-\eta$. By the Kummer description of
paths on $\mathbf G_m$, this implies
\[
 \kappa\left(\frac{\alpha-\eta}{\zeta\alpha-\eta}\right)=0
 \quad\hbox{on }G_{M\Q_p(\mu_\infty)}.
\]
None of these arguments is zero, since $t\ne1$. Formula (C.7b.1)
shows that every displayed ratio is a root of unity.

Embed the finitely generated characteristic-zero field
$\Q(\alpha,\mu_m)$ into $\C$. This does not require $\alpha$ to be
algebraic over $\Q$: if it is transcendental, choose a transcendental
complex image and extend over the finite algebraic part. Roots of unity
have complex absolute value one. Thus the images of $\alpha$ and
$\zeta\alpha$ have equal distances to every point of $\mu_m$. Since
$m\ge3$, those points include three noncollinear centres. Equal distances
to three noncollinear points determine a point of the plane uniquely:
subtracting two squared-distance identities gives two independent real
linear equations for their difference. Hence $\alpha=\zeta\alpha$,
so $\zeta=1$. We have proved zero winding at $0$ modulo every $m\ge3$,
and hence zero winding in $\widehat\Z$. Apply the same argument after
$x\mapsto1-x$ to obtain zero winding at $1$ as well.

For a general genus-zero $V/K$, the rational point identifies its proper
compactification with $\mathbf P^1_K$. Enlarge $K$ finitely to split
all cusps; the given fixed element remains fixed after restriction.
Choose one cusp as infinity. For each other cusp, retain it, infinity,
and one further cusp, and fill in all remaining cusps. A fractional linear
coordinate identifies the resulting curve with the tripod, and its point
remains ordinary. The result just proved kills the selected winding
coordinates. These coordinates form the full geometric abelianization,
which proves the first assertion.

We apply this first assertion again to the covers of the explicit
tripod. Write
\[
 \widehat\Lambda=\widehat\Z[[\widehat\Z]]
   =\varprojlim_{r,m}(\Z/r)[\Z/m],
\]
where the transition systems are ordered by divisibility. The full
profinite free-group universal property gives
\[
 \mathcal N^{\rm ab}\simeq\widehat\Lambda\,b,\qquad
 \mathcal Q\simeq\widehat\Lambda\rtimes\widehat\Z.            \tag{C.7b.2}
\]
Here the generator $T$ of the completed group ring acts by conjugation
by $a$. The proof is identical to the explicit two-way construction in
(C.7a.4), now in the category of all profinite groups: $\mathcal N$ is
the closed normal subgroup generated by $b$, its abelianization is a
cyclic continuous $\widehat\Lambda$-module, and the map of the free
profinite group to $\widehat\Lambda\rtimes\widehat\Z$ sending $b$ to
$1$ supplies the inverse. In the degree-$m$ cover, the image of a module
element in $\widehat\Lambda/(T^m-1)$ consists of its winding coordinates
at the $m$ roots of unity. This follows either by lifting $a^i b a^{-i}$,
which winds once around the corresponding root, or by the free geometric
presentation of the punctured projective line.

For the ordinary point $2$, the first assertion puts a full fixed element
$g$ in $\mathcal N$, hence in every degree-$m$ covering subgroup. Choose
a lift $\alpha^m=2$ and a finite local extension containing it and all
cusps. The lifted loop is fixed after restriction to the corresponding
cyclotomic extension. The first assertion, applied now to the entire
punctured projective line $V_m$ and its ordinary point $\alpha$, makes
all its full winding coordinates zero. Therefore the image of $g$ in
$\widehat\Lambda/(T^m-1)$ is zero for every $m\ge3$. Reduction also
modulo every positive integer $r$ gives all the defining finite quotients
of the completed group ring, so its image in $\widehat\Lambda$ is zero.
This proves the ordinary-point formula.

For the tangent, the full cusp inertia is fixed by $H$. Subtract from a
fixed element $g$ its full cusp-$0$ winding by multiplying by the
corresponding inverse element of $I_0$. The resulting fixed element $h$
has zero $a$-coordinate and belongs to every degree-$m$ power cover.
The unit tangent at $z=0$ lifts the original unit tangent, as is seen from
$x=z^m$ in the Puiseux fibre functor. Fill in $0$ on this cover. The
section becomes the ordinary point $0$ on
$\mathbf P^1\setminus(\mu_m\cup\{\infty\})$, a hyperbolic genus-zero
curve for $m\ge2$. Its full fixed loop has all winding coordinates at
$\mu_m$ zero by the first assertion. The remaining winding $u_0$ on the
unfilled cover satisfies $m u_0=0$ by pushforward to the original
$a$-coordinate; multiplication by $m$ is injective on $\widehat\Z$,
so $u_0=0$. Formula (C.7b.2) now shows that the image of $h$ vanishes,
as in the ordinary case. Thus the tangent image is contained in $I_0$;
that inertia supplies the reverse inclusion.
\end{proof}
